\chapter{Vaping in pregnancy}
\index{Vaping in pregnancy}

\section*{Definition and Overview}
Vaping in pregnancy refers to the use of electronic cigarettes (e-cigarettes) or other vaping devices while pregnant. E-cigarettes are battery-powered devices that heat a liquid containing nicotine, flavouring, and other chemicals to produce an aerosol that is inhaled. They were introduced as a less harmful alternative to smoking, and some people use them to reduce cigarette use or to quit smoking. However, "less harmful than cigarettes" does not mean "harmless," and the safety of vaping during pregnancy has not been established.

Because pregnancy is a period of profound developmental vulnerability, the safest approach is to avoid all nicotine and all inhaled aerosol products. No e-cigarette has been approved as a smoking-cessation aid in pregnancy, and the known risks of nicotine to the developing fetus mean that health authorities, including the local government, recommend that pregnant women do not vape. For women who smoke, the priority is to quit completely, using the safest available support, and to obtain accurate advice rather than silently switching to an unregulated substitute.

\section*{Epidemiology}
Smoking during pregnancy has declined substantially in many countries over recent decades, but it has not disappeared, and vaping has introduced a new and rapidly evolving pattern of use. The local National Drug Strategy Household Survey shows that e-cigarette use has risen sharply, particularly among young people, and some of this growth has occurred in women of reproductive age. A proportion of women who vape continue to do so when they become pregnant, and mixed use of both cigarettes and e-cigarettes is also seen.

Accurate national data on vaping specifically during pregnancy are limited because the practice is newer than cigarette smoking and patterns change quickly. Nevertheless, surveys suggest that a small but meaningful minority of pregnant women report current or recent vaping, with higher rates among younger women, women who smoke, and women who used vaping to quit smoking. The popularity of disposable, fruit- and candy-flavoured devices with high nicotine concentrations has increased overall exposure, and health professionals now routinely ask about vaping alongside smoking as part of antenatal care.

\section*{Aetiology and Pathophysiology}
\subsection*{What is in the aerosol}
E-cigarette aerosol is not water vapour. It contains fine particles and a variable mixture of chemicals, the most concerning of which for pregnancy are nicotine, flavourings, and the metals and organic compounds generated by heating.

\begin{itemize}
    \item \textbf{Nicotine:} a powerful, addictive substance that crosses the placenta freely and is detected in fetal blood at concentrations similar to or higher than maternal levels
    \item \textbf{Flavourings:} chemicals that may be harmless in food but, when inhaled, can irritate airways and some of which are known to damage lung tissue in animal studies
    \item \textbf{Metals:} traces of lead, nickel, chromium, and tin have been found in e-cigarette aerosol, released from the heating coil
    \item \textbf{Carbonyls:} formaldehyde and other compounds can form during heating, especially at high power settings
\end{itemize}

\subsection*{How nicotine affects the fetus}
Nicotine constricts the mother's blood vessels, reducing blood flow through the placenta and therefore the supply of oxygen and nutrients to the baby. In animal studies it has also been shown to act directly on the developing brain, interfering with the formation of synapses, the connections between nerve cells. These effects are relevant to growth restriction, to the baby's brain development, and to the later risk of sudden infant death. Nicotine exposure during pregnancy is a recognised cause of low birth weight regardless of how it is delivered, whether through smoking, vaping, or nicotine replacement products, so the assumption that vaping removes all risk is incorrect.

\subsection*{Comparison with smoking}
Cigarette smoke contains thousands of chemicals, including carbon monoxide, tar, and numerous carcinogens, and carries a far larger burden of harm than e-cigarette aerosol. This is why switching completely from smoking to vaping is considered to reduce harm for a non-pregnant adult who cannot quit in any other way. But during pregnancy the question is different: there is no benefit of nicotine at all, so the safest option is to avoid it entirely rather than to substitute one source for another.

\section*{Clinical Features}
Vaping has no "symptoms" of its own in pregnancy; rather, the concern is the silent effect on the growing baby.

\begin{itemize}
    \item \textbf{Maternal effects:} vaping may cause throat and airway irritation, cough, and headaches; severe lung injury from vaping is rare but has been reported overseas
    \item \textbf{Effects on the baby:} nicotine-related reduction in placental blood flow can lead to slower growth, lower birth weight, and a small increase in the risk of premature birth
    \item \textbf{Brain development:} the developing fetal brain is exposed to nicotine at a time of rapid growth, with animal evidence of long-term effects on attention and behaviour
    \item \textbf{Addiction:} nicotine dependence in the mother can persist after birth, increasing the risk that the baby is exposed to second-hand aerosol in the home
    \item \textbf{Dual use:} many people who vape continue to smoke as well, so any benefit over smoking is lost, and the baby is exposed to both
\end{itemize}

The absence of immediate symptoms is itself part of the risk. A woman may vape throughout pregnancy feeling perfectly well, and the baby may appear healthy at birth, yet the effects of nicotine on growth and neurodevelopment are subtle and measurable across a whole population, not necessarily visible in any individual baby. This is why the advice is unconditional: there is no known safe level of nicotine in pregnancy, and there is no way to tell whether a particular pregnancy is being affected. These risks apply to nicotine replacement products as well, which is why they are used in pregnancy only as a short-term aid to achieve the real goal, complete cessation of all nicotine, rather than as a long-term substitute.

\section*{Differential Diagnosis}
The clinical issue is not a disease to be distinguished from others but a behaviour to be understood and addressed. Relevant considerations include the following.

\begin{itemize}
    \item Exclusive vaping versus continued cigarette smoking or mixed (dual) use
    \item Nicotine-containing versus nicotine-free e-liquid
    \item Vaping used as a quit attempt versus recreational use
    \item Use of other substances, including alcohol, cannabis, or non-prescribed medicines
    \item Underlying respiratory conditions, such as asthma, which vaping can worsen
\end{itemize}

\section*{When to Seek Medical Care}
Every pregnant woman should disclose smoking or vaping at her first antenatal appointment, and ideally at every subsequent visit, because support is most effective when offered early. Seek advice if you vape and are pregnant, thinking about becoming pregnant, or trying to stop. It is also important to seek help if you are finding it difficult to quit on your own, as counselling, behavioural support, and specialist smoking-cessation services significantly improve the chance of success.

\textbf{Contact your local emergency services for:}
\begin{itemize}
    \item Difficulty breathing, chest pain, or coughing up blood after vaping
    \item Sudden severe abdominal pain or reduced fetal movements
    \item Heavy vaginal bleeding or signs of labour
    \item Any sudden, severe, or concerning symptom in pregnancy
\end{itemize}

\section*{Diagnosis and Investigations}
\begin{itemize}
    \item \textbf{Antenatal history:} routine enquiry about smoking and vaping, including products, frequency, and nicotine content, recorded at booking and follow-up visits
    \item \textbf{Carbon monoxide monitoring:} a simple breath test offered to women who smoke or vape, which gives immediate feedback and helps track progress
    \item \textbf{Ultrasound and growth scans:} increased surveillance is offered if there are concerns about fetal growth, so that slower growth can be detected early
    \item \textbf{Referral:} referral to a smoking-cessation service or Quitline for intensive behavioural support, which is the single most effective intervention
\end{itemize}

\section*{Management}
\subsection*{The message}
The clear recommendation from local and international bodies is that pregnant women should not use e-cigarettes. If a woman does not smoke, she should not start vaping in pregnancy. If she smokes, the priority is complete cessation using the most effective and safest supports available, rather than switching to vaping.

\subsection*{Stopping smoking with support}
\begin{itemize}
    \item \textbf{Behavioural support:} counselling through Quitline, midwives, and general practitioners roughly doubles the chance of quitting successfully
    \item \textbf{Nicotine replacement therapy:} patches, gum, or lozenges may be considered under professional supervision when the benefits of quitting outweigh the risks of nicotine exposure; they deliver nicotine without the thousands of toxicants in smoke
    \item \textbf{Prescribed medicines:} bupropion can be used in pregnancy under specialist advice
    \item \textbf{Relapse prevention:} most quit attempts face setbacks, so a practical, non-judgemental plan for managing cravings is essential
\end{itemize}

\subsection*{If the woman chooses to vape}
A clinician's role is to support reduction and eventual cessation rather than to abandon the discussion. Advice focuses on eliminating dual use with cigarettes, reducing nicotine strength and frequency, avoiding use around others, and planning to stop as early as possible. Referral to cessation services remains the central plank of management.

\section*{Complications}
\begin{itemize}
    \item Low birth weight and slower fetal growth related to nicotine's effect on placental blood flow
    \item A small increase in the risk of premature birth
    \item Potential effects on the developing fetal brain, with animal evidence linking nicotine exposure to later difficulties with attention and impulse control
    \item Increased risk of sudden infant death if the baby is exposed to nicotine after birth and to second-hand aerosol in the home
    \item Maternal nicotine dependence that persists after pregnancy
    \item Unknown long-term effects of inhaled flavourings and other aerosol chemicals, which have not been studied over decades of use
\end{itemize}

\section*{Prognosis and Outcomes}
The outlook depends almost entirely on the pattern of use. Continued cigarette smoking throughout pregnancy carries well-documented risks of growth restriction, preterm birth, low birth weight, and infant mortality. Complete cessation early in pregnancy reduces these risks substantially, and women who stop smoking by the end of the first trimester approach the outcomes of non-smokers for most complications. Because it is nicotine rather than the smoke alone that harms fetal growth, vaping during pregnancy is not a risk-free alternative, and the best outcomes are seen in women who achieve complete abstinence from both cigarettes and e-cigarettes. With the right support, most women who smoke or vape can quit during pregnancy, and doing so is one of the most powerful steps a woman can take for the health of her baby.

\section*{Prevention and Self-Care}
\begin{itemize}
    \item If you are planning a pregnancy, stop smoking and vaping before you conceive; this is the ideal time to quit
    \item Tell your midwife, GP, or obstetrician about any vaping so you can get tailored support
    \item Use the free support of Quitline or an antenatal smoking-cessation service rather than trying to go it alone
    \item Do not vape in the home, car, or around children; second-hand aerosol is not harmless
    \item Ask partners and household members to support a smoke-free and vape-free home
    \item Remember that reducing, then stopping, nicotine entirely gives the best chance of a healthy pregnancy
\end{itemize}

\section*{Sources and Further Reading}
The following organisations provide reliable information on vaping and smoking in pregnancy.

\begin{itemize}
    \item NHS. \textit{Information on smoking and vaping in pregnancy and stopping support.} https://www.nhs.uk/conditions/
    \item Mayo Clinic. \textit{Guide to quitting smoking during pregnancy and e-cigarette safety.} https://www.mayoclinic.org/
    \item MedlinePlus. \textit{Health information on smoking and vaping during pregnancy.} https://medlineplus.gov/
    \item World Health Organization. \textit{Resources on tobacco and electronic nicotine products in pregnancy.} https://www.who.int/
    \item Quitline Australia. \textit{Free support for quitting smoking and vaping.} https://www.quit.org.au/
    \item HealthDirect Australia. \textit{Consumer information on vaping and pregnancy.} https://www.healthdirect.gov.au/
\end{itemize}
